\chapter{Plantar Fasciitis}

\section*{Definition}
Plantar fasciitis is a common cause of heel pain, characterized by sharp, stabbing discomfort at the plantar aspect of the heel that is most intense with the first steps in the morning or after prolonged sitting. It involves microtears and degeneration of the plantar fascia at its calcaneal attachment.

\section*{Pathophysiology}
The plantar fascia is a thick, fibrous aponeurosis that supports the longitudinal arch of the foot. Repetitive tensile overload, especially during weight-bearing activities, causes microtrauma at the medial calcaneal tubercle. This triggers a degenerative rather than primarily inflammatory response (fasciosis), with collagen necrosis, angiofibroblastic hyperplasia, chondroid metaplasia, and matrix degradation. Perifascial edema and increased vascularity may be seen in the acute phase.

\section*{Causes}
\begin{itemize}
  \item \textbf{Biomechanical:} Pes planus (flat feet), pes cavus (high arches), overpronation, leg length discrepancy
  \item \textbf{Occupational:} Prolonged standing or walking on hard surfaces
  \item \textbf{Activity-related:} Running, jumping, ballet, aerobic dance; sudden increase in training intensity or duration
  \item \textbf{Systemic:} Obesity, diabetes mellitus, rheumatoid arthritis, reactive arthritis, ankylosing spondylitis
  \item \textbf{Anatomic:} Tight Achilles tendon and gastrocnemius complex, weak intrinsic foot muscles
  \item \textbf{Age:} Most common between 40--60 years
\end{itemize}

\section*{When to See a Doctor}
See your doctor if:
\begin{itemize}
  \item Morning first-step pain lasting more than 4 weeks despite home measures
  \item Severe heel pain limiting walking or daily activities
  \item Associated with numbness, tingling, or burning (suspect tarsal tunnel syndrome)
  \item Bilateral heel pain with systemic symptoms (suspect inflammatory arthritis)
\end{itemize}

\section*{Related Symptoms}
\begin{itemize}
  \item Inferior heel pain, worse in the morning or after rest
  \item Pain with prolonged standing or after weight-bearing activity
  \item Tenderness at the medial calcaneal tubercle on palpation
  \item Pain worsened by dorsiflexion of the toes (windlass test)
  \item Secondary Achilles tendon tightness
\end{itemize}
\textbf{Important:} Bilateral plantar fasciitis with morning stiffness lasting > 30 minutes may indicate an underlying seronegative spondyloarthropathy rather than a mechanical problem.

\section*{Self-Care / Home Management}
\begin{itemize}
  \item Gastrocnemius and plantar fascia stretching before getting out of bed
  \item Arch supports or gel heel cups in shoes
  \item Ice pack rolled under the arch for 10--15 minutes
  \item Over-the-counter NSAIDs (ibuprofen, naproxen) for pain
  \item Avoid walking barefoot on hard surfaces; wear supportive shoes at all times
\end{itemize}

\section*{How Your Doctor Will Evaluate You}
\begin{itemize}
  \item Clinical diagnosis based on characteristic history and physical examination
  \item Ultrasound: plantar fascia thickening > 4 mm, hypoechoic areas, and hyperemia
  \item MRI for atypical features (fascial rupture, fat pad atrophy, occult stress fracture)
  \item Radiography to rule out calcaneal stress fracture, bone spur (incidental), or other osseous pathology
  \item Laboratory evaluation: ESR, CRP, HLA-B27 if inflammatory arthritis is suspected
\end{itemize}

\section*{Treatment Options}
\begin{itemize}
  \item Stretching protocol (gastrocnemius, hamstring, plantar fascia) as primary treatment
  \item NSAIDs for 2--4 weeks for symptomatic relief
  \item Custom orthotics or prefabricated arch supports
  \item Night splints to maintain ankle dorsiflexion during sleep
  \item Corticosteroid injection for short-term relief (avoid multiple injections due to fascial rupture risk)
  \item Extracorporeal shock wave therapy for chronic cases lasting > 6 months
  \item Plantar fascia release (open or endoscopic) for refractory cases after 9--12 months of conservative therapy
\end{itemize}

\clearpage
